% Source: source/普通高考/2022/2022全国乙理(河南,山西,江西,安徽,陕西,甘肃,青海,内蒙古,黑龙江,吉林,.宁夏,新疆).pdf, page 1, question 6.
% PDF embedded-bitmap check: pdfimages -list found no image objects; the flowchart is vector artwork.
% Grid evidence: tmp/redraw_flowchart_2022_national_paper_b_science/pages/page-1_grid100.png,
%   tmp/redraw_flowchart_2022_national_paper_b_science/flow_big_grid50.png;
% comparison original side is cropped from the no-grid page render.
% Elements: start/end rounded boxes, input/output parallelograms, two process boxes,
%   one decision diamond, downward arrows, and the left loop-back branch.
% Node-edge list from the source PDF: 开始 -> 输入; 输入 -> b=b+2a;
%   b=b+2a -> a=b-a,n=n+1; a=b-a,n=n+1 -> decision;
%   decision --是--> 输出 n -> 结束; decision --否--> b=b+2a.
% Relations: the loop returns from the decision's 否 branch to the b=b+2a process;
%   the 是 branch continues to 输出 n and then 结束. The decision tests
%   |b^2/a^2-2|<0.01 exactly as printed in the source.
% solid segments: all box outlines, branch lines, and arrow shafts/heads.
% dashed segments: none.
% labels: Chinese labels and formulas remain centered in their boxes; 否 is left of
%   the diamond branch and 是 is right of the downward branch, clear of the lines.

\begin{tikzpicture}[baseline=(current bounding box.center), y=1.13cm,
  flowtext/.style={anchor=center, align=center, font=\normalsize,
    execute at begin node={\setlength{\parindent}{0pt}}},
  flowbox/.style={flowtext, draw, line width=0.45pt, minimum height=0.48cm,
    inner xsep=0.15cm, inner ysep=0.05cm},
  flowarrow/.style={-Stealth, line width=0.45pt},
]
  % Source figure is a schematic flowchart; coordinates preserve topology, not scale data.
  \node[flowbox, rounded corners=0.12cm, minimum width=0.95cm, minimum height=0.60cm] (start) at (0,0) {开始};
  \coordinate (inputTopL) at (-2.20,-0.73);
  \coordinate (inputTopR) at (2.20,-0.73);
  \coordinate (inputBotR) at (2.00,-1.25);
  \coordinate (inputBotL) at (-2.40,-1.25);
  \node[flowtext] (input) at (0,-0.99) {输入 $a=1,b=1,n=1$};
  \draw[flowbox] (inputTopL) -- (inputTopR) -- (inputBotR) -- (inputBotL) -- cycle;
  \node[flowbox, minimum width=1.88cm, minimum height=0.53cm] (updateb) at (0,-1.95)
    {$b=b+2a$};
  \node[flowbox, minimum width=3.35cm, minimum height=0.55cm] (updatea) at (0,-2.90)
    {$a=b-a,n=n+1$};
  \node[flowtext, draw, line width=0.45pt, diamond, aspect=3.6, minimum width=6.60cm,
    minimum height=1.18cm, inner sep=0pt] (test) at (0,-4.36)
    {$\displaystyle\left|\frac{b^2}{a^2}-2\right|<0.01$};
  \coordinate (outputTopL) at (-0.80,-5.54);
  \coordinate (outputTopR) at (0.80,-5.54);
  \coordinate (outputBotR) at (0.60,-6.06);
  \coordinate (outputBotL) at (-1.00,-6.06);
  \node[flowtext] (output) at (0,-5.83) {输出 $n$};
  \draw[flowbox] (outputTopL) -- (outputTopR) -- (outputBotR) -- (outputBotL) -- cycle;
  \node[flowbox, rounded corners=0.12cm, minimum width=0.95cm, minimum height=0.60cm] (finish) at (0,-6.79) {结束};

  \draw[flowarrow] (start) -- (0,-0.73);
  \draw[flowarrow] (0,-1.25) -- (updateb.north);
  \draw[flowarrow] (updateb) -- (updatea);
  \draw[flowarrow] (updatea) -- (test);
  \draw[flowarrow] (0,-6.12) -- (finish.north);

  % The 否 branch turns left, rises to the level below the input box, and
  % re-enters the b-update process through its top-center arrow.
  \coordinate (loopleft) at (-3.90,-4.28);
  \coordinate (looplevel) at (-3.90,-1.52);
  \draw[line width=0.45pt] (test.west) -- node[flowtext, below] {否}
    (loopleft) -- (looplevel) -- (0,-1.52);
  \draw[flowarrow] (0,-1.52) -- (updateb.north);
  \draw[flowarrow] (test.south) -- node[pos=0.05, flowtext, right] {是} (0,-5.54);
\end{tikzpicture}
